\section{Introduction}\label{sec:intro}

Randomness is often described in one of two ways. On one view, it is a primitive feature of the world: outcomes are genuinely indeterminate. On the other, it is a deficit on the side of the observer: the world is orderly, but the observer lacks the information or computational power required to predict it. Both descriptions leave out an operational fact that scientific theories cannot avoid. A theory never receives the substrate for free. It packages some distinctions into stable objects, ignores others, and then calls the predictive remainder noise, chance, or stochasticity. Without specifying the layer, the packaging, and the budget that maintains it, claims about randomness remain underdetermined.

Within the Six Birds program, closure is organized by six primitive operations: \emph{P1 operator rewrite}, \emph{P2 constraints}, \emph{P3 protocol holonomy / route mismatch}, \emph{P4 staging}, \emph{P5 packaging}, and \emph{P6 accounting} \citep{tsiokos2026sixbirds}. In the present paper, P5 and P6 do most of the work. Packaging decides which substrate distinctions become objects at the layer, while accounting decides which such distinctions can actually be afforded, stabilized, and audited. P4 sets the scale and timescale at which closure is attempted, P2 specifies which distinctions are admissible, P3 captures variability induced by untracked route dependence, and P1 describes the interventions required when closure fails to descend. Randomness is not a seventh primitive. It is the variability that remains when these closure mechanics do not fully stabilize the distinctions that matter for prediction at the chosen layer.

This paper develops that claim in a form that can be computed and tested. Our thesis is that randomness is the predictive residue of non-closure. Let $X_t$ denote a micro process and let $\Pi : X \to Y$ be a packaging map, with $Y_t = \Pi(X_t)$ the packaged state observed at the layer. We define the \emph{micro-closure deficit} at scale $\tau$ by
\[
\mathsf{CD}_\tau(\Pi) := I(X_t; Y_{t+\tau} \mid Y_t).
\]
This quantity measures exactly how much information about the next packaged state is still hidden inside the current macro-object. It yields the decomposition
\[
H(Y_{t+\tau}\mid Y_t)
=
H(Y_{t+\tau}\mid X_t)
+
\mathsf{CD}_\tau(\Pi),
\]
which separates intrinsic substrate uncertainty from randomness introduced by discarded distinctions. A macrostate looks random exactly to the extent that it has failed to become a sufficient statistic for its own future. When the substrate is deterministic, the closure deficit is the whole story. When the substrate itself is stochastic, observable layer-relative randomness decomposes into intrinsic noise and non-closure.

This viewpoint extends and sharpens several earlier strands of the Six Birds sequence. \emph{Six Birds: Foundations of Emergence Calculus} defined packaging as an idempotent form of closure and distinguished it from novelty and directionality \citep{tsiokos2026sixbirds}. \emph{To Notch a Stone with Six Birds} operationalized predictive non-closure through an order-1 versus order-2 log-loss gap, using that gap to motivate theory birth and lens refinement \citep{tsiokos2026notch}. \emph{To Create a Stone with Six Birds} supplied an explicit coarse-graining pipeline for induced macro kernels and route-mismatch diagnostics on stochastic substrates \citep{tsiokos2026pica}. If \emph{To Notch a Stone} asked when failure of closure forces a richer theory to be born, the present paper asks what that failure looks like inside the theory that still tries to run. Our answer is: unresolved closure is experienced as randomness.

The paper also places this Six Birds account in contact with existing external traditions. The compression--prediction trade-off studied by the information bottleneck and related representation-learning programs already treats predictive quality as something purchased under a representational budget \citep{tishby1999bottleneck,geiger2015optimalkullbackleibleraggregation}. Computational mechanics similarly isolates predictive sufficiency as the criterion that separates structure from randomness in a process \citep{shalizi2001computationalmechanics}. Our contribution is different in emphasis. We treat predictive residue explicitly as a \emph{closure deficit} of a packaging, connect it to concrete coarse-graining diagnostics already native to Six Birds, and then carry the same logic into engineered one-wayness. In particular, hashing will appear not as a source of mysterious randomness, but as an especially sharp case of many-to-one packaging under feasibility constraints, in the same spirit that the random-oracle tradition uses idealized random functions to reason about one-way behavior \citep{bellare1993randomoracles}.

Three controlled demonstrations drive the paper. First, on finite Markov chains with known partitions, closure deficit vanishes on exactly lumpable constructions and rises under controlled within-fiber heterogeneity, while route mismatch tracks it closely. This gives a direct bridge between the exact quantity introduced here and the computable coarse-graining diagnostics already used in Six Birds and in classical lumpability analyses of Markov chains \citep{kemeny1976finitemarkovchains,buchholz1994exactordinarylumpability,tsiokos2026pica}. Second, in budget-limited prediction, increasing memory order buys back hidden distinctions that the current packaged state failed to stabilize, producing a concrete randomness--budget curve. Third, in a toy-safe hashing setting, inversion success follows the familiar $q/2^n$ scaling for high-entropy inputs at small success probabilities, while low-entropy inputs collapse the effect entirely. One-wayness and random-lookingness then appear as the same packaging-and-budget phenomenon: not entropy creation, but feasible irreversibility relative to bounded observers.

Our claims are deliberately narrower than the familiar metaphysical debate between ontic and epistemic randomness. We do not claim that all randomness is reducible to hidden microstructure, nor do we identify route mismatch with genuine irreversibility. We claim something operationally sharper: once a substrate, a packaging, and a maintenance budget are fixed, the non-intrinsic part of layer-relative randomness is measurable as the information about the next packaged state that remains hidden inside the current one. That claim is exact, falsifiable, and portable across domains.

All experiments, figure-generation code, and supporting formal files for this manuscript are tracked in the companion repository at \url{https://github.com/ioannist/six-birds-randomness}.

The rest of the paper proceeds as follows. \Cref{sec:setup} briefly recalls the Six Birds setting needed here. \Cref{sec:closure} defines closure deficit and states the exact decomposition. \Cref{sec:diagnostics} connects the exact quantity to route mismatch and budget-limited predictors. \Cref{sec:results-markov,sec:results-budget,sec:results-hashing} present the three controlled demonstrations. \Cref{sec:discussion} discusses scope, limitations, and the distinction between randomness, novelty, and arrow-of-time; \Cref{sec:conclusion} closes.
